% Project entries use one title/date row because they have no separate role line.
\newcommand*{\cvproject}[3]{%
  \vspace{2mm}
  \setlength\tabcolsep{0pt}%
  \setlength{\extrarowheight}{0pt}%
  \begin{tabular*}{\textwidth}{@{\extracolsep{\fill}} L{\textwidth - 4.5cm} R{4.5cm}}%
    \entrytitlestyle{#1} & \entrydatestyle{#2} \\
    \multicolumn{2}{L{\textwidth}}{\descriptionstyle{\vspace{0mm}#3}}%
  \end{tabular*}%
}

%-------------------------------------------------------------------------------
%  SECTION TITLE
%-------------------------------------------------------------------------------
\vspace{4mm}
\cvsection{Personal Projects}

%-------------------------------------------------------------------------------
%  CONTENT
%-------------------------------------------------------------------------------
\begin{cventries}
  \cvproject
  {Facial Image Alignment and Visualisation Tool}
  {2026}
  {
    \begin{cvitems}
      \item {Developed a Python computer vision tool, inspired by my own Invisalign journey, to align facial images and generate consistent progress visualisations.}
      \item {Used OpenCV and MediaPipe for facial landmark detection, image alignment, morph video generation, and side-by-side comparison outputs.}
      \item {Enhanced robustness with alignment quality checks, error handling, and automated output generation.}
    \end{cvitems}
  }

  \cvproject
  {Mortgage Calculator Web Application}
  {2024}
  {
    \begin{cvitems}
      \item {Developed a Python and Streamlit web application to calculate mortgage repayments and affordability based on borrowing assumptions and desired monthly repayments.}
      \item {Used Bank of England yield curve data to support repayment modelling, affordability estimates, and Stamp Duty calculations for first-time buyers and standard residential purchases.}
      \item {Enhanced robustness with automated data updates, improved calculation accuracy, refactoring, test coverage, and clearer validation/error handling.}
    \end{cvitems}
  }

  \cvproject
  {BBC Good Food Scraper and Web Application}
  {2023}
  {
    \begin{cvitems}
      \item {Developed an end-to-end Python data project to improve recipe filtering based on macro-nutrient targets and available ingredients.}
      \item {Used Python, SQL, GitHub Actions, and AWS RDS to automate weekly collection and storage of recipe data in a cloud-hosted database.}
      \item {Created a Streamlit web application to give non-technical users an interactive way to search and filter recipes.}
    \end{cvitems}
  }
\end{cventries}
